% GENERATED FILE. Edit canonical lessons, then run npm run generate:books.
% canonical-source-hash: fnv1a64:d8a30e631050a9f4
% canonical-lessons: KA-C273-haru, KA-C273-suttadu, KA-C273-datu, KA-C273-odisu, KA-C273-jaru

\chapter{To fly, To wander around, To cross, To drive, To slip}
\label{ch:ka-to-fly-to-wander-around-to-cross-to-drive-to-slip}
\hlchaptermodality{273}

\begin{chapteropening}
\textbf{By the end of this chapter:} \emph{I can say to fly, to wander around, to cross, to drive and to slip.}

Complete the last lesson of chapter 273: I can say to fly, to wander around, to cross, to drive and to slip.
\end{chapteropening}

\section[\texorpdfstring{\kn{ಹಾರು} (hāru)}{ಹಾರು (hāru)}]{\kn{ಹಾರು} (hāru) — to fly}

\label{lesson:KA-C273-haru}



\begin{quote}
Before the new one: say the Kannada for to stir, then the Kannada for to roast.
\end{quote}

\subsection*{You'll want to know: \kn{ಹಾರು}}
\textbf{\kn{ಹಾರು}} — \emph{hāru} — \textquotedblleft{}to fly\textquotedblright{}.

\begin{grammarlens}[title={What you've built}]
One more word for this chapter.
\end{grammarlens}

\subsection*{Your turn}
\begin{itemize}
\raggedright
  \item \emph{Say it:} \emph{hāru}
  \item \emph{Say it:} \emph{hāru}, once more
  \item \emph{Say it:} say \emph{hāru}
  \item \emph{Recall:} say the Kannada for to stir, then the Kannada for to roast, then say \emph{hāru} again
\end{itemize}

\begin{tcolorbox}[breakable,colback=teal!4,colframe=teal!35!black,title={Before you move on}]
What does \kn{ಹಾರು} mean? (\textquotedblleft{}To fly\textquotedblright{}.) Say it once more.
\end{tcolorbox}
\section[\texorpdfstring{\kn{ಸುತ್ತಾಡು} (suttāḍu)}{ಸುತ್ತಾಡು (suttāḍu)}]{\kn{ಸುತ್ತಾಡು} (suttāḍu) — to wander around}

\label{lesson:KA-C273-suttadu}



\begin{quote}
Before the new one: say the Kannada for to roast, then the Kannada for to fly.
\end{quote}

\subsection*{You'll want to know: \kn{ಸುತ್ತಾಡು}}
\textbf{\kn{ಸುತ್ತಾಡು}} — \emph{suttāḍu} — \textquotedblleft{}to wander around\textquotedblright{}.

\begin{grammarlens}[title={What you've built}]
One more word for this chapter.
\end{grammarlens}

\subsection*{Your turn}
\begin{itemize}
\raggedright
  \item \emph{Say it:} \emph{suttāḍu}
  \item \emph{Say it:} \emph{suttāḍu}, once more
  \item \emph{Say it:} say \emph{suttāḍu}
  \item \emph{Recall:} say the Kannada for to roast, then the Kannada for to fly, then say \emph{suttāḍu} again
\end{itemize}

\begin{tcolorbox}[breakable,colback=teal!4,colframe=teal!35!black,title={Before you move on}]
What does \kn{ಸುತ್ತಾಡು} mean? (\textquotedblleft{}To wander around\textquotedblright{}.) Say it once more.
\end{tcolorbox}
\section[\texorpdfstring{\kn{ದಾಟು} (dāṭu)}{ದಾಟು (dāṭu)}]{\kn{ದಾಟು} (dāṭu) — to cross}

\label{lesson:KA-C273-datu}



\begin{quote}
Before the new one: say the Kannada for to fly, then the Kannada for to wander around.
\end{quote}

\subsection*{You'll want to know: \kn{ದಾಟು}}
\textbf{\kn{ದಾಟು}} — \emph{dāṭu} — \textquotedblleft{}to cross\textquotedblright{}.

\begin{grammarlens}[title={What you've built}]
One more word for this chapter.
\end{grammarlens}

\subsection*{Your turn}
\begin{itemize}
\raggedright
  \item \emph{Say it:} \emph{dāṭu}
  \item \emph{Say it:} \emph{dāṭu}, once more
  \item \emph{Say it:} say \emph{dāṭu}
  \item \emph{Recall:} say the Kannada for to fly, then the Kannada for to wander around, then say \emph{dāṭu} again
\end{itemize}

\begin{tcolorbox}[breakable,colback=teal!4,colframe=teal!35!black,title={Before you move on}]
What does \kn{ದಾಟು} mean? (\textquotedblleft{}To cross\textquotedblright{}.) Say it once more.
\end{tcolorbox}
\section[\texorpdfstring{\kn{ಓಡಿಸು} (ōḍisu)}{ಓಡಿಸು (ōḍisu)}]{\kn{ಓಡಿಸು} (ōḍisu) — to drive, to ride}

\label{lesson:KA-C273-odisu}



\begin{quote}
Before the new one: say the Kannada for to wander around, then the Kannada for to cross.
\end{quote}

\subsection*{You'll want to know: \kn{ಓಡಿಸು}}
\textbf{\kn{ಓಡಿಸು}} — \emph{ōḍisu} — \textquotedblleft{}to drive, to ride\textquotedblright{}.

\begin{grammarlens}[title={What you've built}]
One more word for this chapter.
\end{grammarlens}

\subsection*{Your turn}
\begin{itemize}
\raggedright
  \item \emph{Say it:} \emph{ōḍisu}
  \item \emph{Say it:} \emph{ōḍisu}, once more
  \item \emph{Say it:} say \emph{ōḍisu}
  \item \emph{Recall:} say the Kannada for to wander around, then the Kannada for to cross, then say \emph{ōḍisu} again
\end{itemize}

\begin{tcolorbox}[breakable,colback=teal!4,colframe=teal!35!black,title={Before you move on}]
What does \kn{ಓಡಿಸು} mean? (\textquotedblleft{}To drive, to ride\textquotedblright{}.) Say it once more.
\end{tcolorbox}
\section[\texorpdfstring{\kn{ಜಾರು} (jāru)}{ಜಾರು (jāru)}]{\kn{ಜಾರು} (jāru) — to slip}

\label{lesson:KA-C273-jaru}



\begin{quote}
Before the new one: say the Kannada for to cross, then the Kannada for to drive.
\end{quote}

\subsection*{You'll want to know: \kn{ಜಾರು}}
\textbf{\kn{ಜಾರು}} — \emph{jāru} — \textquotedblleft{}to slip\textquotedblright{}.

\begin{grammarlens}[title={What you've built}]
One more word for this chapter.
\end{grammarlens}

\subsection*{Your turn}
\begin{itemize}
\raggedright
  \item \emph{Say it:} \emph{jāru}
  \item \emph{Say it:} \emph{jāru}, once more
  \item \emph{Say it:} say \emph{jāru}
  \item \emph{Recall:} say the Kannada for to cross, then the Kannada for to drive, then say \emph{jāru} again
\end{itemize}

\begin{tcolorbox}[breakable,colback=teal!4,colframe=teal!35!black,title={Before you move on}]
What does \kn{ಜಾರು} mean? (\textquotedblleft{}To slip\textquotedblright{}.) Say it once more.
\end{tcolorbox}
